% GENERATED FILE. Edit canonical lessons, then run npm run generate:books.
% canonical-source-hash: fnv1a64:8070921e0b7e4fdf
% canonical-lessons: IT-C281-vigilia, IT-C281-indomani, IT-C281-stamattina, IT-C281-stanotte, IT-C281-mattinata

\chapter{Il giorno prima, il giorno dopo}
\label{ch:it-il-giorno-prima-il-giorno-dopo}
\hlchaptermodality{281}

\begin{chapteropening}
\textbf{By the end of this chapter:} \emph{I can place events in time: the eve, the next day, this morning, last night and the whole morning.}

Complete the last lesson of chapter 281: I can place events in time: the eve, the next day, this morning, last night and the whole morning.
\end{chapteropening}

\section[la vigilia]{la vigilia — the eve, the day before}

\label{lesson:IT-C281-vigilia}



\begin{quote}
Before the new one: say the Italian for lately, then the Italian for previously.
\end{quote}

\subsection*{You'll want to know: la vigilia}
\textbf{la vigilia} — \textquotedblleft{}the eve, the day before\textquotedblright{}.

\begin{quote}
La vigilia dell'esame non ho dormito.
\end{quote}

\begin{quote}
\emph{On the eve of the exam I did not sleep.}
\end{quote}

Above all the eve of a holiday or of a big day.

\begin{cousinweb}
\emph{vigilia} is a learned borrowing from Latin \emph{vigilia} \textquotedblleft{}a staying awake, a night watch\textquotedblright{} (Lewis \& Short), from \emph{vigil} \textquotedblleft{}awake\textquotedblright{}: the night people stayed up before a feast.
\end{cousinweb}

\begin{grammarlens}[title={What you've built}]
One more word for telling, explaining and describing at B1.
\end{grammarlens}

\subsection*{Your turn}
\begin{itemize}
\raggedright
  \item \emph{Say it:} \emph{la vigilia}
  \item \emph{Say it:} \emph{la vigilia}, once more
  \item \emph{Say it:} a sentence of your own with \emph{la vigilia}
  \item \emph{Recall:} say the Italian for lately, then the Italian for previously, then say \emph{la vigilia} again
\end{itemize}

\begin{tcolorbox}[breakable,colback=teal!4,colframe=teal!35!black,title={Before you move on}]
What does la vigilia mean? (\textquotedblleft{}The eve, the day before\textquotedblright{}.) Say it once more.
\end{tcolorbox}
\section[l'indomani]{l'indomani — the next day, the day after}

\label{lesson:IT-C281-indomani}



\begin{quote}
Before the new one: say the Italian for previously, then the Italian for the eve.
\end{quote}

\subsection*{You'll want to know: l'indomani}
\textbf{l'indomani} — \textquotedblleft{}the next day, the day after\textquotedblright{}.

\begin{quote}
L'indomani siamo partiti presto.
\end{quote}

\begin{quote}
\emph{The next day we left early.}
\end{quote}

Narrative and written; in speech, \emph{il giorno dopo}.

\begin{cousinweb}
\emph{indomani} is \emph{in} with \emph{domani} \textquotedblleft{}tomorrow\textquotedblright{}, which is inherited from Late Latin \emph{dē māne} \textquotedblleft{}from the morning\textquotedblright{}; Latin \emph{māne} meant \textquotedblleft{}in the morning\textquotedblright{} (Lewis \& Short).
\end{cousinweb}

\begin{grammarlens}[title={What you've built}]
One more word for telling, explaining and describing at B1.
\end{grammarlens}

\subsection*{Your turn}
\begin{itemize}
\raggedright
  \item \emph{Say it:} \emph{l'indomani}
  \item \emph{Say it:} \emph{l'indomani}, once more
  \item \emph{Say it:} a sentence of your own with \emph{l'indomani}
  \item \emph{Recall:} say the Italian for previously, then the Italian for the eve, then say \emph{l'indomani} again
\end{itemize}

\begin{tcolorbox}[breakable,colback=teal!4,colframe=teal!35!black,title={Before you move on}]
What does l'indomani mean? (\textquotedblleft{}The next day, the day after\textquotedblright{}.) Say it once more.
\end{tcolorbox}
\section[stamattina]{stamattina — this morning}

\label{lesson:IT-C281-stamattina}



\begin{quote}
Before the new one: say the Italian for the eve, then the Italian for the next day.
\end{quote}

\subsection*{You'll want to know: stamattina}
\textbf{stamattina} — \textquotedblleft{}this morning\textquotedblright{}.

\begin{quote}
Stamattina ho perso l'autobus.
\end{quote}

\begin{quote}
\emph{This morning I missed the bus.}
\end{quote}

Short for \emph{questa mattina}; in Tuscany also \emph{stamani}.

\begin{cousinweb}
\emph{stamattina} fuses \emph{questa} \textquotedblleft{}this\textquotedblright{} and \emph{mattina} \textquotedblleft{}morning\textquotedblright{}, which goes back to Latin \emph{mātūtīnus} \textquotedblleft{}of the early morning\textquotedblright{} (Lewis \& Short).
\end{cousinweb}

\begin{grammarlens}[title={What you've built}]
One more word for telling, explaining and describing at B1.
\end{grammarlens}

\subsection*{Your turn}
\begin{itemize}
\raggedright
  \item \emph{Say it:} \emph{stamattina}
  \item \emph{Say it:} \emph{stamattina}, once more
  \item \emph{Say it:} a sentence of your own with \emph{stamattina}
  \item \emph{Recall:} say the Italian for the eve, then the Italian for the next day, then say \emph{stamattina} again
\end{itemize}

\begin{tcolorbox}[breakable,colback=teal!4,colframe=teal!35!black,title={Before you move on}]
What does stamattina mean? (\textquotedblleft{}This morning\textquotedblright{}.) Say it once more.
\end{tcolorbox}
\section[stanotte]{stanotte — last night; tonight}

\label{lesson:IT-C281-stanotte}



\begin{quote}
Before the new one: say the Italian for the next day, then the Italian for this morning.
\end{quote}

\subsection*{You'll want to know: stanotte}
\textbf{stanotte} — \textquotedblleft{}last night; tonight\textquotedblright{}.

\begin{quote}
Stanotte ha nevicato.
\end{quote}

\begin{quote}
\emph{It snowed last night.}
\end{quote}

With a past tense it is last night; with the present or the future, tonight.

\begin{cousinweb}
\emph{stanotte} fuses \emph{questa} \textquotedblleft{}this\textquotedblright{} and \emph{notte} \textquotedblleft{}night\textquotedblright{}, which is inherited from Latin \emph{nox, noctis} \textquotedblleft{}night\textquotedblright{} (Lewis \& Short).
\end{cousinweb}

\begin{grammarlens}[title={What you've built}]
One more word for telling, explaining and describing at B1.
\end{grammarlens}

\subsection*{Your turn}
\begin{itemize}
\raggedright
  \item \emph{Say it:} \emph{stanotte}
  \item \emph{Say it:} \emph{stanotte}, once more
  \item \emph{Say it:} a sentence of your own with \emph{stanotte}
  \item \emph{Recall:} say the Italian for the next day, then the Italian for this morning, then say \emph{stanotte} again
\end{itemize}

\begin{tcolorbox}[breakable,colback=teal!4,colframe=teal!35!black,title={Before you move on}]
What does stanotte mean? (\textquotedblleft{}Last night; tonight\textquotedblright{}.) Say it once more.
\end{tcolorbox}
\section[la mattinata]{la mattinata — the whole morning, the morning as a span}

\label{lesson:IT-C281-mattinata}



\begin{quote}
Before the new one: say the Italian for this morning, then the Italian for last night.
\end{quote}

\subsection*{You'll want to know: la mattinata}
\textbf{la mattinata} — \textquotedblleft{}the whole morning, the morning as a span\textquotedblright{}.

\begin{quote}
Ho passato la mattinata in biblioteca.
\end{quote}

\begin{quote}
\emph{I spent the morning in the library.}
\end{quote}

The ending \emph{-ata} stresses the span of time, as in \emph{la giornata}.

\begin{cousinweb}
\emph{mattinata} is \emph{mattina} \textquotedblleft{}morning\textquotedblright{} with \emph{-ata}, the ending for a stretch of time; \emph{mattina} goes back to Latin \emph{mātūtīnus} \textquotedblleft{}of the early morning\textquotedblright{} (Lewis \& Short).
\end{cousinweb}

\begin{grammarlens}[title={What you've built}]
One more word for telling, explaining and describing at B1.
\end{grammarlens}

\subsection*{Your turn}
\begin{itemize}
\raggedright
  \item \emph{Say it:} \emph{la mattinata}
  \item \emph{Say it:} \emph{la mattinata}, once more
  \item \emph{Say it:} a sentence of your own with \emph{la mattinata}
  \item \emph{Recall:} say the Italian for this morning, then the Italian for last night, then say \emph{la mattinata} again
\end{itemize}

\begin{tcolorbox}[breakable,colback=teal!4,colframe=teal!35!black,title={Before you move on}]
What does la mattinata mean? (\textquotedblleft{}The whole morning, the morning as a span\textquotedblright{}.) Say it once more.
\end{tcolorbox}
